% TODO make hw stylesheet
\documentclass[10pt]{article}
\usepackage[letterpaper,top=1in]{geometry}
\usepackage[dvipsnames,svgnames]{xcolor}
\usepackage[shortlabels]{enumitem}
\usepackage[framemethod=TikZ]{mdframed}
\usepackage{amsmath,amssymb,amsthm}
\usepackage[colorlinks]{hyperref}
\usepackage{multicol}
\usepackage{tabularx}

\hypersetup{
	colorlinks=true,
	linkcolor=blue,
	filecolor=magenta,
	urlcolor=magenta,
	pdftitle={MAT 324 HW}
}

\usepackage{mathtools}
\usepackage{thmtools}
\usepackage[textsize=footnotesize]{todonotes}
\usepackage{titling}

\reversemarginpar
\mdfdefinestyle{mdbluebox}{roundcorner=10pt,innerbottommargin=9pt,
	linecolor=blue,backgroundcolor=TealBlue!5,}
\declaretheoremstyle[headfont=\sffamily\bfseries\color{MidnightBlue},
mdframed={style=mdbluebox},]{thmbluebox}

\mdfdefinestyle{mdbloobox}{frametitlefont=\bfseries,innerbottommargin=8pt,
	nobreak=true,backgroundcolor=TealBlue!5,linecolor=blue,}
\declaretheoremstyle[headfont=\bfseries\color{MidnightBlue},
mdframed={style=mdbloobox},headpunct={\\[3pt]},postheadspace=0pt,]{thmbloobox}

\mdfdefinestyle{mdredbox}{frametitlefont=\bfseries,innerbottommargin=8pt,
	nobreak=true,backgroundcolor=Salmon!5,linecolor=RawSienna,}
\declaretheoremstyle[headfont=\bfseries\color{RawSienna},
mdframed={style=mdredbox},headpunct={\\[3pt]},postheadspace=0pt,]{thmredbox}

\mdfdefinestyle{mdgreenbox}{linecolor=ForestGreen,backgroundcolor=ForestGreen!5,
	linewidth=2pt,rightline=false,leftline=true,topline=false,bottomline=false,}
\declaretheoremstyle[headfont=\bfseries\sffamily\color{ForestGreen!70!black},
mdframed={style=mdgreenbox},headpunct={ --- },]{thmgreenbox}

\mdfdefinestyle{mdorangebox}{linecolor=RedOrange,backgroundcolor=RedOrange!5,
	linewidth=2pt,rightline=false,leftline=true,topline=false,bottomline=false,}
\declaretheoremstyle[headfont=\bfseries\sffamily\color{RedOrange!70!black},
mdframed={style=mdorangebox},headpunct={ --- },]{thmorangebox}

\mdfdefinestyle{mdblackbox}{linecolor=black,backgroundcolor=RedViolet!5!gray!5,
	linewidth=3pt,rightline=false,leftline=true,topline=false,bottomline=false,}
\declaretheoremstyle[mdframed={style=mdblackbox}]{thmblackbox}

\declaretheoremstyle[spaceabove=6pt,spacebelow=6pt]{basehead}

\declaretheorem[style=basehead,name=Problem]{problem}
\declaretheorem[style=thmredbox,name=Problem,numbered=no]{problem*}
\declaretheorem[style=thmbloobox,name=Setup]{setup} % for config geo
\declaretheorem[style=thmredbox,name=Example,sibling=problem]{example}
\declaretheorem[style=thmredbox,name=$\bigstar$ Problem,sibling=problem]{niceprob}
\declaretheorem[style=thmbluebox,name=Theorem,sibling=problem]{theorem}
\declaretheorem[style=thmbluebox,name=Corollary,sibling=theorem]{corollary}
\declaretheorem[style=thmbluebox,name=Proposition,sibling=theorem]{prop}
\declaretheorem[style=thmbluebox,name=Lemma,numberwithin=problem]{lemma}
\declaretheorem[style=thmbluebox,name=Theorem,numbered=no]{theorem*}
\declaretheorem[style=thmbluebox,name=Lemma,numbered=no]{lemma*}
\declaretheorem[style=thmgreenbox,name=Claim,numberwithin=problem]{claim}
\declaretheorem[style=thmgreenbox,name=Claim,numbered=no]{claim*}
\declaretheorem[style=thmblackbox,name=Remark,numberwithin=problem]{remark}
\declaretheorem[style=thmblackbox,name=Remark,numbered=no]{remark*}
\declaretheorem[style=thmgreenbox,name=Definition,sibling=theorem]{definition}
\declaretheorem[style=thmgreenbox,name=Definition,numbered=no]{definition*}

\providecommand{\ol}{\overline}
\providecommand{\ceil}[1]{\left\lceil #1 \right\rceil}
\providecommand{\floor}[1]{\left\lfloor #1 \right\rfloor}
\newcommand{\dd}{\,\mathrm{d}}
\providecommand{\ul}{\underline}
\providecommand{\eps}{\varepsilon}
\providecommand{\half}{\frac{1}{2}}
\providecommand{\CC}{\mathbb C}
\providecommand{\FF}{\mathbb F}
\providecommand{\NN}{\mathbb N}
\providecommand{\QQ}{\mathbb Q}
\providecommand{\RR}{\mathbb R}
\providecommand{\ZZ}{\mathbb Z}
\providecommand{\dg}{^\circ}
\providecommand{\ii}{\item}
\providecommand{\setdiff}{\smallsetminus}
\providecommand{\alert}[1]{{\sffamily\textbf{\textcolor{blue}{#1}}}}
\providecommand{\inv}{^{-1}}
\providecommand{\mb}[1]{\mathbf{#1}}

\newenvironment{soln}{\begin{proof}[Solution]}{\end{proof}}

\providecommand{\pow}{\operatorname{Pow}}
\providecommand{\vocab}[1]{{\textbf{\textcolor{ForestGreen}{#1}}}}
\providecommand{\lcm}{\operatorname{lcm}}
\providecommand{\abs}[1]{\left|#1\right|}

\usepackage{mathrsfs}

% place title at top right
\setlength{\droptitle}{-8ex}
\pretitle{\begin{flushright}\Large\bfseries}
\posttitle{\par\end{flushright}}
\preauthor{\begin{flushright}}
\postauthor{\end{flushright}}
\predate{\begin{flushright}}
\postdate{\end{flushright}}

\title{MAT 324 HW04}
\author{\large Aditya Pahuja}
\date{\large Due 09/25}
\begin{document}
\maketitle

\begin{problem}
    Prove or give a counterexample:
    If $F\subseteq\RR$ is such that every bounded continuous function
    from $F$ to $\RR$ can be extended to a continuous function
    from $\RR$ to $\RR$, then $F$ is a closed subset of $\RR$.
\end{problem}

\begin{soln}
    Suppose $F\subseteq\RR$ is not closed; then, $F$ does not contain
    one of its limit points $p\in\RR$. This means we can find
    a strictly monotone sequence $x_1,x_2,\dots\in F$ converging to $p$;
    suppose without loss of generality that this sequence is increasing.
    Define $g\colon F\to\RR$ by
    \begin{equation*}
        g(x) = \begin{cases}
	    (-1)^n & x = x_n\text{ for some integer }n\\
	    tg(x_n) + (1-t)g(x_{n-1}) &
	    x = tx_n + (1-t)x_{n-1}\text{ for some }t\in[0,1]
	    \text{ and some }n\in\NN\\
	    1 & \text{otherwise}.
        \end{cases}
    \end{equation*}
    This function is continuous on $(-\infty,p)\cap F$ by the pasting lemma because
    it is continuous on the locally finite collection of closed intervals
    $(-\infty,x_1]$ and $[x_n,x_{n+1}]$ for $n\in\NN$;
    the collection is locally finite because
    \begin{enumerate}[(1)]
        \ii it is of course locally finite on the interior of each of these intervals;
	\ii for any $x_n$, the neighborhood $(x_{n-1},x_{n+1})$
	intersects only $[x_{n-1},x_n]$ and $[x_n,x_{n+1}]$
	and no others in the collection.
    \end{enumerate}
    Also, $g$ is continuous on $(p,\infty)\cap F$ because it is constantly $1$
    on $(p,\infty)\cap F$. Therefore, $g$ is continuous on $F$,
    and it is bounded between $-1$ and $1$.

    However, $g$ cannot be extended to a continuous function $\RR\to\RR$
    because for any $\delta>0$, the neighborhood
    $(p-\delta,p+\delta)$ contains the subsequence $x_N,x_{N+1},\dots$
    of $(x_n)_{n\in\NN}$, which implies that $g$ achieves
    both the values $1$ and $-1$ in any open neighborhood of $p$;
    that is, $\lim_{x\to p} g(x)$ doesn't exist, which means that
    any extension $\widetilde g$ of $g$ cannot be continuous at $p$.

    That is, no non-closed subset $F$ of $\RR$ has the property that
    bounded continuous functions on $F$ can be extended to
    continuous functions on $\RR$, so contrapositively
    any set with the extension property must be closed in $\RR$.
\end{soln}

\pagebreak

\begin{problem}
    \begin{enumerate}[(a)]
	\ii Let $A,B_1,B_2,\ldots\subseteq\RR$ and
	$B = \bigcup_{n=1}^\infty \bigcap_{k=n}^\infty B_k$. Show that
	\begin{equation*}
	    B\setminus A\subseteq\bigcup_{k=n}^\infty (B_k\setminus A)
	    \quad\text{and}\quad
	    A\setminus B\subseteq\bigcup_{k=n}^\infty (A\setminus B_k)
	\end{equation*}
	for all $n\in\NN$.
	\ii Show that $A\subseteq\RR$ is Lebesgue if and only if
	for every $\eps>0$ there exists a Borel subset $B\subseteq\RR$
	such that
	\begin{equation*}
	    \abs{(B\setminus A)\cup(A\setminus B)} < \eps.
	\end{equation*}
    \end{enumerate}
\end{problem}

\begin{soln}
    (a) We can write
    \begin{equation*}
	B\setminus A = \bigcup_{n=1}^\infty\left(
	\left(\bigcap_{k=n}^\infty B_k\right)\setminus A\right)
	= \bigcup_{n=1}^\infty \bigcap_{k=n}^\infty (B_k\setminus A)
    \end{equation*}
    Since $\bigcup_{k=n}^\infty (B_k\setminus A)$ is
    an increasing sequence of sets in $n$, we can get rid of
    the first $N$ terms without changing the union; that is,
    \begin{equation*}
	B\setminus A = \bigcup_{n=N}^\infty \bigcap_{k=n}^\infty (B_k\setminus A)
	\subseteq \bigcup_{n=N}^\infty \bigcup_{k=n}^\infty (B_k\setminus A)
	= \bigcup_{k=N}^\infty (B_k\setminus A)
    \end{equation*}
    for any $N\in\NN$. For the other statement we similarly have
    \begin{equation*}
	A\setminus B =
	\bigcap_{n=1}^\infty \bigcup_{k=n}^\infty (A\setminus B_k)
	= \bigcap_{n=N}^\infty \bigcup_{k=n}^\infty (A\setminus B_k)
	\subseteq \bigcup_{n=N}^\infty \bigcup_{k=n}^\infty (A\setminus B_k)
	= \bigcup_{k=N}^\infty (A\setminus B_k).
    \end{equation*}
    for any $N\in\NN$ because $\bigcup_{k=n}^\infty (A\setminus B_k)$ is
    a decreasing sequence in $N$ so we can discard
    the first $N$ terms in the intersection.
    \\

    (b) If $A\subseteq\RR$ is Lebesgue, then Axler 2.71c shows that
    there exists a Borel set $B\subseteq A$ (manifested as the union of
    countably many open subsets of $\RR$) such that $\abs{A\setminus B} = 0$.
    Of course $B\subseteq A$ implies $\abs{B\setminus A} = 0$, so
    \begin{equation*}
	\abs{(B\setminus A)\cup(A\setminus B)} = 0 < \eps\text{ for all }\eps>0.
    \end{equation*}

    Conversely, suppose that $A$ is a subset of $\RR$ such that
    given $\delta>0$, there exists a Borel set $B_\delta$ such that
    $\abs{(A\setminus B_\delta)\cup(B_\delta\setminus A)} < \delta$,
    and for brevity, let
    $D_\delta = (A\setminus B_\delta)\cup(B_\delta\setminus A)$.
    Then, let $\eps>0$. By definition of outer measure, there exists
    an open set $U\subseteq\RR$ such that $D_{\eps/3}\subseteq U$ and
    \begin{equation*}
	\abs U < \eps/3 + \abs{D_{\eps/3}} < 2\eps/3.
    \end{equation*}
    This means $B_{\eps/3}\setminus U$ is a Borel set contained in $A$ such that
    \begin{align*}
	\abs{A\setminus(B_{\eps/3}\setminus U)}
	= \abs{A\cap(B_{\eps/3}\cap U^c)^c}
	&= \abs{A\cap(B_{\eps/3}^c\cup U)}\\
	&= \abs{(A\setminus B_{\eps/3})\cup(A\cap U)}
	\le \abs{D_\delta} + \abs{U}
	< \eps/3 + 2\eps/3 = \eps.
    \end{align*}
    We showed in class that $A$ is Lebesgue if and only if
    for each $\eps>0$ there exists a Lebesgue set $A^-_\eps\subseteq\RR$
    such that $\abs{A\setminus A^-_\eps} < \eps$.
    Here existence is shown by taking $A^-_\eps = B_{\eps/3}\setminus U$,
    which is sufficient to conclude that $A$ is Lebesgue.
\end{soln}

\pagebreak

\begin{problem}
    \begin{enumerate}[(a)]
        \ii Suppose $(X,\mathcal S,\mu)$ is a measure space and
	$f\colon X\to[0,\infty]$ is an $\mathcal S$-measurable function.
	Prove that
	\begin{equation*}
	    \int_X f\dd\mu > 0\text{ if and only if }
	    \mu(\{x\in X : f(x) > 0\}) > 0.
	\end{equation*}
	\ii Give an example of a Borel measurable function
	$f\colon[0,1]\to(0,\infty)$ such that
	the lower Riemann integral $L(f,[0,1])$ is $0$.
	\ii Compute $\int f\dd\lambda$,
	where $f$ is your answer for (b) and
	$\lambda$ is the Lebesgue measure on $[0,1]$.
    \end{enumerate}
\end{problem}

\begin{soln}
    (a) Let $E = \{x\in X : f(x) > 0\}$ and $E_n = \{x\in X : f(x) > 1/n\}$.
    We have $\frac 1n\mathbf 1_{E_n}\le f$, so
    \begin{equation*}
	\frac{\mu(E_n)}n = \int_X \frac 1n\mathbf 1_{E_n}\dd\mu
	\le \int_Xf
    \end{equation*}
    Since $E = \bigcup_{n=1}^\infty E_n$, $\mu(E) = \lim_{n\to\infty}\mu(E_n)$,
    so if $\mu(E) > 0$ there exists $N$ such that $\mu(E_N) > 0$.
    This implies that integral of $f$ is positive because
    it is bounded below by $\frac{\mu(E_N)}{N}$.

    Conversely suppose $\mu(E) = 0$. If $P$ is a $\mathcal S$-partition of $X$,
    then for $I\in P$, or equivalently $\inf_I f \ne 0$ if and only if $I\subseteq E$.
    Therefore
    \begin{equation*}
	\mathcal L(f,P) = \sum_{I\in P} (\inf_If)\mu(I)
	= \sum_{\substack{I\in P\\ I\subseteq E}} (\inf_If)
	\underbrace{\mu(I)}_{\le \mu(E) = 0} = 0
    \end{equation*}
    for every $\mathcal S$-partition of $X$, so taking the supremum gives
    $\int_Xf\dd\mu = 0$.
    \\

    (b) Define $f\colon[0,1]\to(0,\infty)$ by
    \begin{equation*}
        f(x) = \begin{cases}
	    1 & \text{if }x\notin\QQ\\
	    1/q & \text{if }x=p/q\in\QQ
	    \text{ with }p,q\in\ZZ_{\ge0}\text{ and }\gcd(p,q)=1
        \end{cases}
    \end{equation*}
    where $\ZZ_{\ge0}$ is the set of nonnegative integers
    (note that this means $f(0) = 1$).
    Then, any interval $I\subseteq[0,1]$ with more than one point
    contains infinitely many rationals. Since there are only
    finitely many rational numbers with a given denominator,
    $I$ contains rational numbers with arbitrarily large denominator, so
    \begin{equation*}
	\inf_{x\in I} f(x) \le \frac 1N
    \end{equation*}
    for arbitrarily large $N$, i.e. $\inf_{x\in I}f(x) = 0$.
    This means that any interval $I$ in a partition $P$ of $[0,1]$
    contributes $0\cdot\ell(I) = 0$ to the lower Riemann sum $L(f,P)$,
    so the supremum of $L(f,P)$ over all partitions $P$ is $L(f,[0,1]) = 0$.

    It thus remains to verify that $f$ is Borel measurable.
    Since the image of $f$ is the set $\{1/n : n\in\NN\}$,
    for any $S\subseteq(0,\infty)$, the set $f\inv(S)$
    is the union of sets of the form $f\inv(1/n)$ for some $n\in\NN$,
    each of which is Borel since $f\inv(1)$ is $[0,1]\setminus\QQ$,
    which is the complement of the Borel set
    $[0,1]\cap\QQ = \bigcup_{q\in[0,1]\cap\QQ}\{q\}$,
    and $f\inv(1/n)$ for $n>1$ is finite.
    There are only countably many sets of the form $f\inv(1/n)$,
    so in fact $f\inv(S)$ for \emph{any} $S\subseteq(0,\infty)$ is
    a countable union of Borel sets, hence is a Borel set.
    \\

    (c) The integral of $f$ with respect to $\lambda$ is $\boxed1$ because
    \begin{align*}
	1 = \lambda([0,1]\setminus\QQ)
	&= \int_{[0,1]}\mathbf 1_{[0,1]\setminus\QQ}\dd\lambda\\
	&\le \int_{[0,1]} f\dd\lambda
	\le \int_{[0,1]} \mathbf 1_{[0,1]}\dd\lambda = \lambda([0,1]) = 1,
    \end{align*}
    where the inequalities are due to
    $0 < f(x)\le 1$ for rational $x$ and $f(x) = 1$ for irrational $x$.
\end{soln}

\begin{problem}
    Suppose $\mu$ is the measure on $(\NN, 2^\NN)$ defined by
    \begin{equation*}
	\mu(E) = \sum_{n\in E} \frac 1{2^n}.
    \end{equation*}
    Prove that for every $\eps>0$, there exists a set $E\subseteq\NN$
    with $\mu(\NN\setminus E) < \eps$ such that
    any sequence of functions $f_1,f_2,\dots\colon\NN\to\RR$
    that converges pointwise on $\NN$ also converges uniformly on $E$.
\end{problem}

\begin{soln}
    Given $\eps>0$, pick $N$ such that $\frac 1{2^N} < \eps$
    and set $E = \{1, \dots, N\}$, so that
    $\NN\setminus E = \{N+1, N+2, \dots\}$ has measure
    \begin{equation*}
	\sum_{n=N+1}^\infty \frac 1{2^n} = \frac 1{2^N} < \eps.
    \end{equation*}
    Let $\eps'>0$. For each $k=1,\dots,N$, pointwise convergence of the sequence
    $f_1(k)$, $f_2(k)$, \dots means that there exists $M_k>0$ such that
    $n>M_k$ implies $\abs{f(k) - f_n(k)} < \eps'$.
    Take $M = \max_{1\le k\le N} M_k$; then $n>M$ implies
    $\abs{f(k) - f_n(k)} < \eps'$ for all $k\in E$.
    This is what it means for $f_1$, $f_2$, \dots
    to converge uniformly to $f$ on $E$, so we are done.
\end{soln}

\begin{problem}
    Let $(X,\mathcal S,\mu)$ be a measure space.
    \begin{enumerate}[(a)]
	\ii Let $f_n\colon X\to[0,\infty]$ be a sequence of measurable functions.
	Show that
	\begin{equation*}
	    \int_X\left(\sum_{n=1}^\infty f_n\right)\dd\mu =
	    \sum_{n=1}^\infty \left(\int_X f_n\dd\mu\right).
	\end{equation*}
	\ii Let $f\colon X\to[0,\infty)$ be a measurable function and
	\begin{equation*}
	    E_n = \{x\in X : 2^n < f(x) \le 2^{n+1}\}
	\end{equation*}
	for all $n\in\ZZ$. Show that $\int_X f\dd\mu<\infty$
	if and only if $\sum_{n\in\ZZ} 2^n\mu(E_n) < \infty$.
	\ii Let $\alpha\in\RR$ and $\lambda$ be the Lebesgue measure on $(0,1)$.
	Show that $\int_{(0,1)} x^{-\alpha}\dd\lambda<\infty$ if and only if
	$\alpha < 1$.
    \end{enumerate}
\end{problem}

\begin{soln}
    (a) Let $g_N = \sum_{n=1}^N f_n$ and
    $g = \lim_{N\to\infty} g_N = \sum_{n=1}^\infty f_n$.
    Since the $f_n$ are all nonnegative we can see that
    the $g_N$ are a sequence of increasing functions,
    so the monotone convergence theorem yields
    \begin{align*}
	\int_X \left(\sum_{n=1}^\infty f_n\right)\dd\mu
	&= \int_X g\dd\mu
	= \lim_{N\to\infty} \int_X g_N\dd\mu\\
	&= \lim_{N\to\infty} \int_X \left(\sum_{n=1}^N f_n\right)\dd\mu
	= \lim_{N\to\infty} \sum_{n=1}^N\left(\int_X f_n\dd\mu\right)
	= \sum_{n=1}^\infty \left(\int_X f_n\dd\mu\right).
    \end{align*}

    (b) Observe that the $E_n$ are pairwise disjoint and cover $X$,
    and let $g(x) = \sum_{n\in\ZZ} 2^n\cdot\mathbf 1_{E_n}(x)$
    be a simple function. Since $2^n < f(x) \le 2^{n+1}$ for all $x\in E_n$,
    we have $g(x) \le f(x) < 2g(x)$ for all $x\in X$.
    Therefore, if $\int_X f\dd\mu < \infty$, then
    \begin{equation*}
	\sum_{n\in\ZZ}2^n\mu(E_n) =
	\int_X g\dd\mu \le \int_X f\dd\mu < \infty,
    \end{equation*}
    and if conversely the sum is finite, then
    \begin{equation*}
	\int_X f\dd\mu \le \int_X 2g\dd\mu = \sum_{n\in\ZZ}2^{n+1}\mu(E_n)
	= 2\cdot\sum_{n\in\ZZ}2^n\mu(E_n) < \infty.
    \end{equation*}
    \pagebreak

    (c) Given $\alpha\in\RR$ define
    \begin{equation*}
	E_n = \{x\in (0,1) : 2^n < x^{-\alpha} \le 2^{n+1}\}
	= \{x\in (0,1) : 2^{-(n+1)/\alpha} < x \le 2^{-n/\alpha}\}.
    \end{equation*}

    If $\alpha=0$ then $x^{-\alpha} = 1$ for all $x\in(0,1)$
    so the integral is $\int_{(0,1)}1\dd\lambda = 1 < \infty$.

    If $\alpha < 0$, then $-\alpha > 0$ so $0 < x^{-\alpha} < 1$.
    This means that
    \begin{equation*}
	\int_{(0,1)} x^{-\alpha}\dd\lambda < \int_{(0,1)} \dd\lambda
	= \lambda((0,1)) = 1 < \infty.
    \end{equation*}

    If $\alpha > 0$, then
    since $x\in(0,1)$, $2^{-(n+1)/\alpha} < x$ can only happen
    if $n$ is nonnegative; otherwise, if $n\le -1$, then
    $-(n+1)/\alpha \ge 0$ so $2^{-(n+1)/\alpha}\ge 1$.
    For these nonnegative $n$,
    $\lambda(E_n) = 2^{-n/\alpha} - 2^{-(n+1)/\alpha}$.
    The series
    \begin{align*}
	\sum_{n\in\ZZ} 2^n\lambda(E_n) = \sum_{n=1}^\infty 2^n\lambda(E_n)
	&= \sum_{n=1}^\infty 2^n(2^{-n/\alpha} - 2^{-(n+1)/\alpha}) \\
	&= \sum_{n=1}^\infty (2^{n(1-1/\alpha)} - 2^{n(1-1/\alpha)-1/\alpha}) \\
	&= \sum_{n=1}^\infty (1-2^{-1/\alpha})(2^{n(1-1/\alpha)})
    \end{align*}
    is geometric with ratio $2^{1-1/\alpha}$, so this series is finite
    if and only if $\abs{2^{1-1/\alpha}} < 1$, which only holds when
    $1 - \frac 1\alpha < 0$. Since $\alpha$ is positive we can multiply through
    by $\alpha$ and rearrange to get that $\alpha < 1$.
    By part (b), the finiteness of this series is equivalent
    to the finiteness of $\int_{(0,1)}x^{-\alpha}\dd\lambda$,
    so if $0 < \alpha < 1$ then the integral is finite.

    The cases above show that the integral $\int_{(0,1)} x^{-\alpha}\dd\mu$
    is finite if and only if $\alpha < 1$, as desired.
\end{soln}










\end{document}
